\begin{corollary}[arbitrary connected deformation chains]\hypertarget{tag-033C}{}\label{push:src:E.2.4}
\StatementRecord{push:src:E.2.4}{proved}{033C}{\SeriesRef{push:src:E.2.5}, \SeriesRef{push:src:E.2.24}, \SeriesRef{push:src:E.2.21}, \SeriesRef{push:src:E.2.22}, \SeriesRef{push:src:E.2.23}}

The ordered coprime pair $(3k,4l)$ is constant in every
proper smooth holomorphic family over a connected complex
analytic base which contains a constructed sphere member.
Two constructed members are deformation equivalent through
such families, allowing arbitrary intermediate members and
finite chains, if and only if their pairs agree.
Moreover different pairs cannot occur in one continuous
family in the $C^\infty$ space of integrable complex
structures on $S^6$.
\end{corollary}


\begin{proof}

\begin{enumerate}[label={$\langle1\rangle$\arabic*.},leftmargin=*]
\item The detailed proof is \ref{push:src:E.2.5}--\ref{push:src:E.2.24} below.  The intrinsic
pencil persists for either canonical slope, the limiting
period stays in the admissible interior, and the two finite
root charts and the actual analytic cusp are recovered.
\ref{push:src:E.2.21} then identifies every full punctured-disc smooth
completion with the prescribed member.  The additional
step \ref{push:src:E.2.22} propagates from any nonempty known open subset
of a disc: finitely many successive discrete analytic
exceptional sets are filled in reverse order.  \ref{push:src:E.2.23} uses
a finite local projection and normalization of a curve
through two specified points to pass to arbitrary analytic
bases and deformation chains.
The classical $C^\infty$ Kuranishi local theorem then
gives the continuous-family conclusion in \ref{push:src:E.2.24}.

\item Local miniversality alone would prove only openness.
Nor does a boundary point of an arbitrary open locus
necessarily admit a holomorphic disc whose punctured
image lies in that locus.  The generic-disc propagation
and finite-projection steps are the substantive additions
which establish the required closedness.

\end{enumerate}
\end{proof}